% ================================================================
%  conteudo/revisao.tex
%  Indicar apenas os tópicos que deveriam ser abordados.
%
%  Exemplos de citação (ABNT NBR 10520:2023):
%    Autor no texto  → Segundo Xavier (2013), ...
%    Autor ao final  → ... (Xavier, 2013).
%    Dois autores    → (Clarac; Bonnin, 1985)
%    Até 4 autores   → (Cruz; Correa; Costa; Santos, 1998)
%    Mais de 4       → (Urani \textit{et al.}, 1994)
% ================================================================

\section{Revisão da Literatura}

A revisão da literatura deve apresentar os principais conceitos teóricos que embasam a prática realizada. Abaixo estão listados os tópicos que devem ser abordados nesta seção, seguidos de exemplos de como redigir cada um deles com as devidas citações no formato ABNT.

\subsection{Fermentação e metabolismo microbiano}

A fermentação é definida como um processo catabólico no qual compostos orgânicos atuam simultaneamente como doadores e aceptores de elétrons, sem a participação de oxigênio molecular como aceptor final (Lehninger; Nelson; Cox, 2014). Segundo Xavier (2013), processos fermentativos são fundamentais na indústria biotecnológica moderna, sendo empregados desde a produção de etanol até a síntese de antibióticos e enzimas de uso industrial.

Outra forma de citação no final da sentença, quando o(s) autor(es) não está(ão) inserido(s) na frase: a levedura \textit{Saccharomyces cerevisiae} é o microrganismo mais utilizado em fermentações industriais de interesse alimentício e energético (Urani \textit{et al.}, 1994). Dois autores: a regulação do metabolismo fermentativo envolve mecanismos de controle alostérico e repressão catabólica (Clarac; Bonnin, 1985). Até quatro autores: a composição do meio de cultivo influencia diretamente a velocidade e o rendimento do processo (Cruz; Correa; Costa; Santos, 1998).

\subsection{Cinética de crescimento microbiano}

A cinética de crescimento microbiano descreve matematicamente as diferentes fases do desenvolvimento celular em um meio de cultivo. De acordo com Schmidell \textit{et al.} (2001), o crescimento em batelada apresenta quatro fases distintas: (i) fase \textit{lag}, em que os microrganismos se adaptam ao novo meio; (ii) fase exponencial, caracterizada pelo crescimento à velocidade específica máxima ($\mu_{max}$); (iii) fase estacionária, quando o crescimento cessa pela limitação de nutrientes ou acúmulo de metabólitos tóxicos; e (iv) fase de declínio, com morte celular progressiva.

\subsection{Parâmetros cinéticos}

Os principais parâmetros utilizados para caracterizar um processo fermentativo incluem a velocidade específica máxima de crescimento ($\mu_{max}$), o fator de conversão substrato-biomassa ($Y_{X/S}$) e a produtividade volumétrica ($P_x$). Esses parâmetros são essenciais para o dimensionamento e controle de biorreatores em escala industrial (Costa \textit{et al.}, 2020).

